\documentclass[tikz,border=3mm,multi=tikzpicture]{standalone}
\usepackage{eleve-quimica}

\begin{document}

% FIGURA 1 — fig-01-escala-logaritmica-ph
\begin{tikzpicture}[quimica figura, x=1cm, y=1cm]
  \path[use as bounding box] (-0.3,-0.2) rectangle (14.7,6.1);

  \fill[quimicalaranjaclaro] (0.35,2.45) rectangle (6.85,3.45);
  \fill[quimicaverdeclaro] (6.85,2.45) rectangle (7.85,3.45);
  \fill[quimicaazulclaro] (7.85,2.45) rectangle (14.35,3.45);
  \draw[quimicacinza, line width=1.1pt] (0.35,2.45) rectangle (14.35,3.45);

  \foreach \p in {0,...,14}{
    \draw[quimicacinza, line width=0.8pt]
      (0.35+\p,2.45) -- (0.35+\p,2.20);
    \node[quimica rotulo, font=\normalsize\bfseries]
      at (0.35+\p,1.78) {\p};
  }

  \node[quimica destaque] at (3.35,4.05) {meio ácido};
  \node[quimica rotulo, text=quimicaverde] at (7.35,4.05)
    {neutro a 25\,°C};
  \node[quimica rotulo, text=quimicaazul] at (11.35,4.05)
    {meio básico};

  \draw[quimica retorno] (7.00,5.15) -- (0.75,5.15);
  \node[quimica destaque] at (3.85,5.62)
    {$[\mathrm{H_3O^+}]$ aumenta};
  \draw[quimica fluxo] (7.70,5.15) -- (13.95,5.15);
  \node[quimica rotulo, text=quimicaazul] at (10.85,5.62)
    {$[\mathrm{OH^-}]$ aumenta};

  \node[quimica rotulo] at (7.35,0.72)
    {cada unidade: concentração varia por fator 10};
\end{tikzpicture}


% FIGURA 2 — fig-02-curva-titulacao-forte-forte
\begin{tikzpicture}[quimica figura, x=1cm, y=1cm]
  \path[use as bounding box] (-0.2,-0.25) rectangle (12.2,8.55);

  \draw[quimica eixo] (1.10,0.80) -- (11.65,0.80)
    node[above left=2pt, quimica rotulo] {volume de base adicionado};
  \draw[quimica eixo] (1.10,0.80) -- (1.10,8.05)
    node[above=2pt, quimica rotulo] {pH};

  \foreach \p/\y in {0/0.80,2/1.80,4/2.80,6/3.80,7/4.30,8/4.80,10/5.80,12/6.80,14/7.80}{
    \draw[quimicacinza, line width=0.8pt] (0.92,\y) -- (1.10,\y);
    \node[quimica rotulo, font=\normalsize\bfseries, anchor=east]
      at (0.80,\y) {\p};
  }

  \fill[quimicalaranjaclaro, opacity=0.75] (1.10,4.90) rectangle (11.20,5.80);
  \node[quimica destaque, anchor=west] at (7.60,5.38)
    {faixa de viragem};

  \draw[quimica auxiliar] (6.15,0.80) -- (6.15,7.80);
  \draw[quimica auxiliar] (1.10,4.30) -- (6.15,4.30);
  \node[quimica rotulo, anchor=west] at (6.42,4.38)
    {equivalência};
  \node[quimica rotulo, anchor=west] at (6.42,3.95)
    {$\mathrm{pH}=7$};

  \draw[quimica direta]
    (1.35,1.25)
    .. controls (3.00,1.30) and (4.85,1.55) .. (5.70,2.35)
    .. controls (6.00,2.75) and (6.05,3.55) .. (6.15,4.30)
    .. controls (6.25,5.50) and (6.45,6.65) .. (7.15,7.05)
    .. controls (8.20,7.55) and (9.80,7.62) .. (11.15,7.65);
  \node[quimica destaque, anchor=west] at (2.00,1.72)
    {ácido forte};
  \node[quimica rotulo, text=quimicaazul, anchor=east] at (10.95,7.20)
    {excesso de base};
\end{tikzpicture}

\end{document}
